\section{Independent transversals and quota packing}
\label{app:packing}

We give a finite self-contained proof of the graph-theoretic packing theorem
used in Section~\ref{sec:ltar-provider}.

Let $G=(V,E)$ be a finite simple graph of maximum degree $\Delta$, and let
$\mathcal U$ be a finite partition of $V$ into nonempty blocks.  A
\emph{partial independent transversal} is an independent set $S\subseteq V$
meeting each block in at most one point.

For such an $S$, let $\mathcal U_{\rm uncov}(S)$ be the uncovered blocks.  If
$X\subseteq S$, write $\mathcal U(X)$ for the blocks meeting $X$.  Since $S$
meets each block at most once,
\[
|\mathcal U(X)|=|X|.
\tag{H1}
\]

\subsection{Augmenting sequences}

For a sequence $\mathbf v=(v_1,\ldots,v_m)$ define
\[
\mathcal B_0=\mathcal U_{\rm uncov}(S),
\]
\[
\mathcal B_k
=\mathcal B_{k-1}\cup\mathcal U(N(v_k)\cap S),
\]
and
\[
C_0=\varnothing,
\qquad
C_k=C_{k-1}\cup(N(v_k)\cap S)\cup\{v_k\}.
\]
Put
\[
d_k=|N(v_k)\cap S|.
\]
We call $\mathbf v$ augmenting if
\begin{enumerate}
\item $v_k$ lies in a block of $\mathcal B_{k-1}$;
\item $N(v_k)\cap C_{k-1}=\varnothing$;
\item $d_k>0$ whenever $k<m$.
\end{enumerate}
The empty sequence is augmenting.

If $i<j$, then
\[
(N(v_i)\cap S)\cap(N(v_j)\cap S)=\varnothing,
\tag{H2}
\]
because the first set lies in $C_{j-1}$ while the second is disjoint from
$C_{j-1}$.  In particular the vertices $v_i$ are distinct.  Using (H1) and
(H2),
\[
|\mathcal B_m|
=|\mathcal U_{\rm uncov}(S)|+
\sum_{i=1}^m d_i,
\tag{H3}
\]
while
\[
|C_m|\le m+\sum_{i=1}^m d_i.
\tag{H4}
\]

\subsection{Extension while all current degrees are positive}

Assume some block is uncovered and $d_i>0$ for every $i\le m$.  Put
\[
D_m=\sum_{i=1}^m d_i.
\]
Then $m\le D_m$, so by (H3)--(H4),
\[
|C_m|\le2D_m<2|\mathcal B_m|.
\tag{H5}
\]
Suppose every partition block has at least $2\Delta$ vertices.  If
$V_{\mathcal B_m}$ is the union of the blocks in $\mathcal B_m$, then
\[
|V_{\mathcal B_m}|
\ge2\Delta|\mathcal B_m|.
\]
On the other hand
\[
|N(C_m)\cap V_{\mathcal B_m}|
\le\Delta|C_m|
<2\Delta|\mathcal B_m|.
\]
Hence there is
\[
w\in V_{\mathcal B_m}\setminus N(C_m).
\]
Appending $w$ preserves the augmenting conditions.  Thus an augmenting sequence
can be extended whenever uncovered blocks remain and all its current degrees
are positive.  The case $\Delta=0$ is immediate.

\subsection{Improvement when the last degree vanishes}

Suppose $\mathbf v$ is nonempty and $d_m=0$.  Let $U$ be the partition block
containing $v_m$.  Since
\[
N(v_m)\cap S=\varnothing,
\]
there are two cases.

If $U$ is uncovered, then
\[
S'=S\cup\{v_m\}
\]
is a partial independent transversal with $|S'|=|S|+1$.

If $U$ is already covered, let $s$ be the unique point of $S\cap U$.  Since
$U\in\mathcal B_{m-1}$, there is a first $k<m$ with
\[
s\in N(v_k)\cap S.
\]
Replace $s$ by $v_m$:
\[
S'=(S\setminus\{s\})\cup\{v_m\}.
\]
Then $S'$ is again a partial independent transversal, with the same cardinality
and the same covered blocks.  For $j<k$ the quantities
$|N(v_j)\cap S|$ are unchanged, while at $k$ the degree drops by one.  The
prefix $(v_1,\ldots,v_k)$ remains augmenting for $S'$.

Thus a zero last degree yields either a larger partial transversal or, at the
same size, an augmenting sequence whose degree vector decreases at its first
changed coordinate.

\subsection{Well-founded minimization}

Let $M=|V|$.  For a feasible pair
\[
t=(S,(v_1,\ldots,v_m)),
\]
pad the degree vector $(d_1,\ldots,d_m)$ to length $M$ with the sentinel
$M+1$, and order feasible pairs lexicographically by
\[
\mu(t)=
\bigl(M-|S|,\widehat D(t)\bigr).
\]
The range of $\mu$ is finite, so a minimal feasible pair exists.

If its partial transversal were not full, then:
\begin{itemize}
\item if the augmenting sequence is empty, the extension argument applies;
\item if its last degree is zero, the preceding improvement applies;
\item if its last degree is positive, all its degrees are positive by the
  definition of augmenting sequence, so again the extension argument applies.
\end{itemize}
Every case strictly decreases $\mu$, a contradiction.

We have proved the following theorem.

\begin{theorem}[Independent transversal]\label{thm:haxell-internal}
Let a finite graph of maximum degree $\Delta$ be partitioned into nonempty
blocks, each of cardinality at least $2\Delta$.  Then it has an independent
transversal.
\end{theorem}

\subsection{Quota packing}

Let a finite graph of maximum degree at most $d>0$ be partitioned into rows
$R_i$.  In every row split off
\[
\left\lfloor\frac{|R_i|}{2d}\right\rfloor
\]
disjoint chunks of cardinality $2d$, discarding fewer than $2d$ residual
vertices.  Restrict the graph to the union of these chunks; its maximum degree
is still at most $d$.  Theorem~\ref{thm:haxell-internal} supplies one independent
representative from every chunk.  Their union is an independent set meeting
each row in exactly
\[
\left\lfloor\frac{|R_i|}{2d}\right\rfloor
\]
vertices.  For $d=0$, the whole vertex set is independent.

Thus Theorem~\ref{thm:quota-packing} holds with
\[
L(0)=1,
\qquad
L(d)=2d\quad(d>0).
\]
The sharpness of this denominator is not used in the proof of Erdős Problem
289.
